\chapter[Joint detector calibration]{Joint detector calibration and an entropy certificate}
\label{ch:dc-calibration}

Finite detector records need not supply a pointwise volatility estimate in
order to certify dissipation.  Under an explicit structural relation between
the detector coordinate and its local diffusivity, simultaneous moment bounds
can instead calibrate the unknown geometric scale and a polynomial entropy
score.  The construction below accounts for exposure, independent replica
noise and a calibrated detector gain.  It does not infer the structural
relation itself from the records.

\section{Physical and acquisition assumptions}

The hidden process is a stationary smooth uniformly elliptic diffusion on a
smooth compact connected manifold without boundary, with positive smooth
stationary density $\pi$.  Reversal is ordinary hidden-state reversal.  Write
$L=S+A$, $L^\dagger=S-A$, $Af=v\cdot\nabla f$, and
\[
 \Gamma(f,g)=\nabla f^{\mathsf T}D\nabla g,\qquad
 e=v^{\mathsf T}D^{-1}v,\qquad \sigma=\mathbb E_\pi e.
\]
In this convention $d[f(X)]_t=2\Gamma_f(X_t)\,dt$.  For a smooth real
coordinate $f$ with $|f|\le1$, assume the pointwise identity
\begin{equation}
 \Gamma_f(x)=\gamma\bigl(1-f(x)^2\bigr),\qquad
 \gamma>0\text{ unknown and constant}.
 \label{eq:dc-shape}
\end{equation}
Hidden dimension and drift are otherwise unrestricted in the stated class.
An average relation, a stationary marginal or an estimated conditional mean
of $\Gamma_f$ is not a substitute for~\eqref{eq:dc-shape}.  The generator
bounds $M_{1,\pm},M_{2,\pm}$ from the finite-resolution chapter are retained:
\begin{align*}
 \|Lf\|_\infty&\le M_{1,+},&
 \|L^\dagger f\|_\infty&\le M_{1,-},\\
 \|L^2f\|_\infty&\le M_{2,+},&
 \|(L^\dagger)^2f\|_\infty&\le M_{2,-}.
\end{align*}
These are physical calibration assumptions.  Set
$M_*:=\min(M_{1,+},M_{1,-})$ and
$M_{2,*}:=\min(M_{2,+},M_{2,-})$.

Each independent stationary trial has nominal centers $-h,0,h$ and a common
causal uniform exposure
\begin{equation}
 A_t=\frac1\tau\int_{t-\tau}^t f(X_s)\,ds,
 \qquad 0\le\tau\le\tau_U\le h.
 \label{eq:dc-exposure}
\end{equation}
At $\tau=0$, interpret $A_t=f(X_t)$.  Acquire $K=2m+2$ center replicas and
one endpoint readout at each of $-h,h$.  Conditional on the entire hidden
path, all detector randomizations in a trial are independent and each
bounded readout $Z$ has mean $A_t$.  A binary shot with mean $A_t$ is an
example.  Conditional noise independence does not assert independent hidden
states.  Unknown gain is treated separately below.  Repeated powers of a
single noisy readout do not satisfy this replica contract.

\section{An exactly unbiased replica polynomial}

Let $e_k$ be the normalized elementary symmetric polynomial of the $K$
center readouts: the average of all products indexed by distinct
$k$-subsets.  Set $e_0=1$.  The replica assumptions imply
$\mathbb E[e_k\mid X_{\mathbb R}]=A_0^k$.
For fixed $\lambda\in\mathbb R$, put
\begin{align}
 b&=1-\lambda h,&
 \beta_g&=\frac h4(M_{2,+}+M_{2,-}),\\
 \beta_K&=\frac h2(M_{1,-}M_{2,+}+M_{1,+}M_{2,-}),&
 \beta&=2|\lambda|\beta_g+\beta_K.
 \label{eq:dc-beta}
\end{align}
Define the record score
\begin{equation}
 T_m=\frac1{h^2}\sum_{j=0}^m
 \left[-Z_-Z_+e_{2j}
       +b(Z_-+Z_+)e_{2j+1}-b^2e_{2j+2}\right]
       -\beta\sum_{j=0}^m e_{2j}.
 \label{eq:dc-replica-score}
\end{equation}
Its expectation is exactly the displayed polynomial in $A_-,A_0,A_+$.
Write $T_m^{\rm ideal}$ for that polynomial evaluated at
$Y_t=f(X_t)$ instead.

\begin{proposition}[Ideal polynomial inequality]
Under~\eqref{eq:dc-shape},
\begin{equation}
 \mathbb E T_m^{\rm ideal}\le\gamma\sigma.
 \label{eq:dc-polynomial-physical}
\end{equation}
\end{proposition}
\begin{proof}
Set $w_m(y)=\sum_{j=0}^m y^{2j}$ and $q(y)=-\lambda y$.  The uncorrected
ideal polynomial is
\[
 -w_m(Y_0)
 \left(q(Y_0)-\frac{Y_h-Y_0}{h}\right)
 \left(q(Y_0)-\frac{Y_{-h}-Y_0}{h}\right).
\]
The two-sided semigroup bounds in the finite-resolution chapter control its
conditional time-step bias by $2|q|\beta_g+\beta_K\le\beta$.  Conditional on
the complete center state, the corrected expectation is therefore at most
\[
 w_m(f)\{2qSf-q^2-(Sf)^2+(Af)^2\}
 =w_m(f)\{(Af)^2-(q-Sf)^2\}.
\]
Since $(Af)^2\le\Gamma_f e$ and
$w_m(f)\Gamma_f=\gamma(1-f^{2m+2})\le\gamma$, integration proves the claim.
The argument does not clip individual records or intermediate expectations.
\end{proof}

\section{A multitime moment bound for exposure}

\begin{lemma}[Bidirectional gap derivative]
For the stationary ordered moment
$F(t_1,\ldots,t_d)=\mathbb E\prod_i f(X_{t_i})$, its derivative in gap $j$
satisfies
\begin{equation}
 \left|\partial_{t_{j+1}-t_j}F\right|
 \le a_{d,j}:=
 \min\{j^2M_{1,-},(d-j)^2M_{1,+}\},\qquad 1\le j<d.
 \label{eq:dc-gap}
\end{equation}
The assertion includes one-sided limits at coincident times.
\end{lemma}
\begin{proof}
A backward semigroup word $g$ containing $k$ factors $f$ obeys
$\|g\|_\infty\le1$ and $\|Lg\|_{L^1(\pi)}\le k^2M_{1,+}$.
For the induction, semigroup application preserves this bound by stationary
$L^1$ contraction and commutation with $L$.  Multiplication uses
\begin{gather*}
 L(fg)=fLg+gLf+2\Gamma(f,g),\\
 \mathbb E\Gamma_f=-\mathbb E fLf\le M_{1,+},\qquad
 \mathbb E\Gamma_g=-\mathbb E gLg\le\|Lg\|_1.
\end{gather*}
The Cauchy inequality for $\Gamma$ gives the new bound
$B+M_{1,+}+2\sqrt{BM_{1,+}}$ when the old bound is $B$; if
$B\le(k-1)^2M_{1,+}$, the result is at most $k^2M_{1,+}$.
At a gap, the conditional product of preceding factors has absolute value
at most one.  Differentiating the suffix and integrating its generator in
$L^1(\pi)$ yields the $(d-j)^2M_{1,+}$ bound.  Applying the same argument
under stationary time reversal gives $j^2M_{1,-}$.  Smooth compactness
justifies the generator identities and differentiation.
\end{proof}

For arbitrary causal exposure supported in $[t_i-\tau,t_i]$, sorted
displacements from the nominal ordered times lie in $[-\tau,0]$, so each
gap changes by at most $\tau$.  Consequently a degree-$d$ product changes by
at most
\begin{equation}
 \tau\sum_{j=1}^{d-1}a_{d,j}.
 \label{eq:dc-support-blur}
\end{equation}
Products of exposure averages are integrals over auxiliary exposure times;
these integration variables do not imply independence of signal values.

Common uniform boxes give explicit sharper constants.  Suppose the nominal
windows do not overlap and have chronological multiplicities
$(p_1,\ldots,p_r)$, deleting zero entries, with $d=\sum_i p_i$.  Define
$B(p_1,\ldots,p_r)$ by summing the following contributions at each global
prefix $k$.  An internal gap of a cluster with $p$ factors contributes
\begin{equation}
 \frac{a_{d,k}}{p+1}.
 \label{eq:dc-internal-spacing}
\end{equation}
A gap between consecutive clusters with $p$ and $q$ factors contributes
\begin{equation}
 a_{d,k}\left[1-\frac1{p+1}-\frac1{q+1}
                  +\frac{2p!q!}{(p+q+1)!}\right].
 \label{eq:dc-cross-spacing}
\end{equation}
The resulting blur radius is $\tau_U B(p_1,\ldots,p_r)$; degrees zero and
one have radius zero.  To verify the formula, internal uniform order-statistic
spacings have mean $1/(p+1)$.  A cross-gap displacement divided by $\tau$ is
$U_{\min,q}-U_{\max,p}$.  Its positive part has expectation
\[
 \int_0^1t^p(1-t)^q\,dt=\frac{p!q!}{(p+q+1)!},
\]
while the mean of its negative is
$1-1/(p+1)-1/(q+1)$.  This yields its expected absolute value in
\eqref{eq:dc-cross-spacing}.  Integrate~\eqref{eq:dc-gap} along the segment
between the two nonnegative gap vectors and then average.

\section{Calibrating the geometric scale}

Let
\[
 C(t)=\mathbb E[Y_0Y_t],\quad D_0=\mathbb E A_0^2,\quad
 D_h=\mathbb E[A_0A_h],\quad q_0=\mathbb E\Gamma_f.
\]
Stationarity implies $C(-t)=C(t)$, $C'(0+)=-q_0$, $q_0\le M_*$ and
$|C''(t)|\le M_{2,*}$ for $t\ge0$, with a right derivative at zero.  Hence
\begin{align}
 0\le C(0)-D_0&\le\delta_0:=\frac{M_*\tau_U}{3},\\
 |D_h-C(h)|&\le\delta_h:=\frac{M_{2,*}\tau_U^2}{12}.
 \label{eq:dc-covariance-blur}
\end{align}
The first bound uses $C(t)\le C(0)$, the Lipschitz constant $M_*$ and
$\mathbb E|U-V|=\tau/3$ for two uniform exposure offsets.  The second uses
their centered difference, whose variance is $\tau^2/6$, in the Taylor
remainder; $h\ge\tau$ prevents passage through the covariance cusp.  The
same bound holds at $2h$.  A more precise identity is
$D_0=C(0)-q_0\tau/3+R_0$, with $|R_0|\le M_{2,*}\tau^2/12$; the displayed
bound avoids an unknown $q_0$ in the calibration envelope.

Write $\mathcal D(p_1,\ldots,p_r)$ for either valid general blur radius
above, optionally replacing a degree-two radius by $\delta_0$ for one window
and by $\delta_h$ for distinct windows separated by at least $h$.  Define
\begin{align}
 \delta_N={}&\frac1{h^2}\sum_{j=0}^m
 \bigl\{\mathcal D(1,2j,1)
 +|b|[\mathcal D(1,2j+1)+\mathcal D(2j+1,1)]
 +b^2\mathcal D(2j+2)\bigr\}\\
 &\quad+\beta\sum_{j=0}^m\mathcal D(2j).
 \label{eq:dc-score-blur}
\end{align}
Delete zero clusters before evaluation.  The last term accounts for blur of
the bias-correction moments.  The two cross terms need not have identical
bounds when the forward and adjoint caps differ.

The following three population quantities now have compatible directions:
\begin{align}
 P&=\mathbb ET_m-\delta_N\le\gamma\sigma,\\
 \overline A&=1-D_0-\delta_0\le A_{\rm true}:=1-C(0),\\
 \overline H&=\frac{D_0-D_h+\delta_0+\delta_h}{h}
                     +\frac{M_{2,*}h}{2}
                 \ge\gamma A_{\rm true}.
 \label{eq:dc-three-population}
\end{align}
The last inequality follows from
$|C(0)-C(h)-hq_0|\le M_{2,*}h^2/2$ and the exact relation
$q_0=\gamma(1-C(0))$.  This step uses neither an arcsine marginal nor a
local-volatility estimate.

\section{A simultaneous finite-record certificate}

\begin{theorem}[Joint geometry and score calibration]
Assume simultaneous bounds $p\le P$, $a\le\overline A$ and
$u\ge\overline H$ hold with probability at least $1-\alpha$ for the
declared acquisition.  Set
\begin{equation}
 \widehat L=
 \begin{cases}
 [p]_+[a]_+/u,&u>0,\\
 0,&u\le0,
 \end{cases}
 \label{eq:dc-certificate}
\end{equation}
marking the latter case inconclusive.  Then
$\Pr(\sigma\ge\widehat L)\ge1-\alpha$.  On the same confidence event, if
$a>0$, then $\gamma\le u/a$.
\end{theorem}
\begin{proof}
If $p,a>0$, then $A_{\rm true}>0$ and
$u\ge\gamma A_{\rm true}>0$.  Therefore
$pa/u\le\sigma a/A_{\rm true}\le\sigma$.  Otherwise the displayed lower
bound is zero.  The geometric upper bound follows from
$u\ge\gamma A_{\rm true}\ge\gamma a$.
\end{proof}

One construction covers the means of $T_m$, $e_2$ and
$(Z_-+Z_+)e_1/2$ simultaneously using the same independent trials.  Their
population means are respectively $\mathbb ET_m,D_0,D_h$.  Interval
arithmetic supplies the required lower and upper endpoints.  Unconditional
bounded-mean concentration for each statistic, followed by a declared union
bound, remains valid when the statistics share records.  There is no need
to split data merely because the geometric scale and the entropy score are
calibrated together.

For $|Z|\le B$, $|e_k|\le B^k$, so valid finite score ranges follow directly.
Hoeffding or a valid empirical-Bernstein inequality supplies the mean
intervals.  Numerical endpoints require outward enclosure.  A variance
estimate cannot be substituted into a theorem requiring a known variance.
Selecting $m$ or $\lambda$ from the same records is valid only if the
underlying moment region or a finite-family confidence event covers that
selection.  Applying an ordinary iid score inequality after an uncovered
coefficient selection would not establish the theorem's premise.

\section{Unknown detector gain and a physical reference protocol}

Suppose a raw bounded readout $R$ has mean $\kappa A_t$ conditional on the
hidden path, with unknown common constant $\kappa>0$.  A raw degree-$d$
product has mean $\kappa^d$ times its exposed moment.  Cover its raw mean and
project that interval over a simultaneously covered positive gain region;
this is different from inserting an estimated gain as if exact.
In the formulas of this gain-calibration section and the raw-contrast
projection below, $e_k$ denotes the elementary polynomial of the raw center
readouts $R$, before gain correction.

For example, the complete raw contrast
\begin{equation}
 S_j=\frac1{h^2}\left[-R_-R_+e_{2j}
       +b(R_-+R_+)e_{2j+1}-b^2e_{2j+2}\right]
 \label{eq:dc-gain-contrast}
\end{equation}
has homogeneous degree $2j+2$.  Its mean interval can be divided by
$\kappa^{2j+2}$ while retaining cancellations within the contrast.  The
degree-$2j$ bias-correction means are handled separately.  Signed endpoints
must be minimized or maximized over the gain interval: dividing every
endpoint by its lower gain is invalid for negative means.  A full joint
region may preserve the common gain across degrees; rectangular projection
is conservative.  If the gain region reaches zero, finite inverse bounds
may be unavailable and the procedure can abstain.

A concrete bounded reference protocol can calibrate both gain and exposure.
Feed the same detector a known constant signal one, whose mean readout is
$\kappa$.  Also feed a known step from zero to one at time zero.  If a
coarse metrological bracket permits a fixed reference time
$0<s<\tau_{\min}\le\tau$, the step's mean readout is
\begin{equation}
 p_{\rm step}=\kappa s/\tau.
 \label{eq:dc-step}
\end{equation}
Simultaneous bounded-mean confidence intervals yield the joint constraints
\[
 \kappa\in[\kappa_L,\kappa_U],\qquad
 p_{\rm step}\in[p_L,p_U],\qquad p_{\rm step}\tau=\kappa s.
\]
When $p_L>0$, a conservative exposure interval is
\begin{equation}
 \left[\frac{\kappa_Ls}{p_U},\frac{\kappa_Us}{p_L}\right]
 \label{eq:dc-shutter-region}
\end{equation}
intersected with the coarse bracket.  Without the ramp condition the correct
response is $\kappa\min(1,s/\tau)$, and that relation must be retained.
An empty region is an inconclusive/model-check flag.  A nonpositive lower
step mean can leave the exposure upper bound inconclusive.

These references require a known input shape, a known time origin and common
detector response between reference and target.  They do not validate
\eqref{eq:dc-shape} or the generator caps.  Cover all raw target and reference
means with one declared error allocation.  Dependence among statistics from
the same trials is harmless to a union bound; independence across the trials
used by each concentration inequality is still required.  Adaptive reference
design requires its own valid confidence argument.

\section{Retaining covariance before nuisance projection}

Intervals for separate polynomial components can discard useful cancellations.
A predetermined raw contrast can retain those cancellations before concentration,
while still allowing an uncertain gain afterward.  Choose a nominal
$\kappa_0>0$ as a predeclared fixed design constant, or using a pilot
independent of the inference records, and include in the same covered mean
vector, all $S_j$, the required $e_{2j}$, the calibration contrast, and
\begin{equation}
 R_\ell=\sum_{j=0}^{\ell}\frac{S_j}{\kappa_0^{2j+2}},
 \qquad 0\le\ell\le m.
 \label{eq:dc-predetermined-contrast}
\end{equation}
Here the $S_j$ include the $h^{-2}$ factor in
\eqref{eq:dc-gain-contrast}.  A confidence interval for $\mathbb ER_\ell$
is computed from its actual record variance or range, not by summing the
individual $S_j$ margins.  This is the established principle of concentrating
a fixed linear contrast before projecting nuisance uncertainty.

Suppose $\mathbb ES_j\in[L_j,U_j]$ and
$\mathbb ER_\ell\ge L_{R_\ell}$ on the common confidence event.  For
$d_j=2j+2$, define
\[
 I_j=[\kappa_U^{-d_j}-\kappa_0^{-d_j},
       \kappa_L^{-d_j}-\kappa_0^{-d_j}],\qquad
 c_j=\min\{xy:x\in\partial I_j,\ y\in\{L_j,U_j\}\}.
\]
The four-corner minimum is valid because a bilinear function on a rectangle
attains its extrema at corners.  For every gain in the covered interval,
\begin{equation}
 \sum_{j=0}^{\ell}\frac{\mathbb ES_j}{\kappa^{2j+2}}
 =\mathbb ER_\ell+
   \sum_{j=0}^{\ell}(\kappa^{-d_j}-\kappa_0^{-d_j})\mathbb ES_j
 \ge L_{R_\ell}+\sum_{j=0}^{\ell}c_j.
 \label{eq:dc-correlation-projection}
\end{equation}
Subtract the covered upper bound on
$\beta\sum_{j=0}^{\ell}\mathbb E e_{2j}/\kappa^{2j}$ and the appropriate
$\delta_N$ to obtain a lower bound for $P_\ell$.  The maximum of this bound
and the componentwise lower bound remains valid on the same event.  All
prefixes $\ell\le m$ use the same $K=2m+2$ genuine replicas: their elementary
averages remain unbiased.  Once the common confidence vector covers every
contrast actually used, choosing the largest certificate across these
prefixes spends no further error probability.  Adding an uncovered contrast
after inspecting the data would require additional justification.

\section{An endpoint entropy comparison}

The order-zero raw contrast also provides a comparison that does not require
the structural identity~\eqref{eq:dc-shape}.  Keep the stationary ordinary
diffusion and detector assumptions and let $0\le b\le1$.  Suppose the common
confidence event gives $s$ below the ideal, gain-corrected and blur-corrected
mean
\begin{equation}
 N_0=\frac{2bC(h)-C(2h)-b^2C(0)}{h^2}.
 \label{eq:dc-endpoint-target}
\end{equation}
This endpoint comparison subtracts the raw order-zero blur radius, but not
the $\beta$ penalty used by the derivative-based entropy score.  Then
\begin{equation}
 \sigma\ge\frac{2h^3[s]_+^2}{(2b+\sqrt2)^2}.
 \label{eq:dc-endpoint-baseline}
\end{equation}
This is the earlier reversible-support and entropy-bridge construction,
expressed in the current covariance normalization.

To check constants, compare $L$ with $S=(L+L^\dagger)/2$ using the same
stationary density, diffusion tensor and coordinate.  The diffusion
Girsanov identity gives
$D(P_L^{[0,t]}\Vert P_S^{[0,t]})=\sigma t/4$.  Pinsker and the bounded
endpoint product yield
$|C_L(t)-C_S(t)|\le\sqrt{\sigma t/2}$.  The zero-lag moments agree exactly,
and reversibility gives
\[
 2bC_S(h)-C_S(2h)-b^2C_S(0)
   =-\|(e^{hS}-bI)f\|_{L^2(\pi)}^2\le0.
\]
It follows that $h^2[N_0]_+\le(2b+\sqrt2)\sqrt{\sigma h/2}$, which proves
\eqref{eq:dc-endpoint-baseline}.  Exposure errors must be removed before
using the ideal zero-lag cancellation.  Both entropy bounds can be reported
on the same simultaneous confidence event.  Their difference then reflects
different physical inequalities, with detector and calibration access held
fixed, rather than a comparison against an uncorrected estimator.

\section[Sharp limiting hierarchy]{A sharp limiting hierarchy with unknown scale}

For the stationary circle
$d\phi=\omega\,dt+\sqrt{2\gamma}\,dW$ and $f=\cos\phi$, one has
$Sf=-\gamma f$ and $(Af)^2=\omega^2(1-f^2)$.  The zero-step weighted target
divided by $\gamma$ is exactly
\begin{equation}
 J_m(\lambda)=\frac{\omega^2}{\gamma}
       \bigl(1-\mathbb EY^{2m+2}\bigr)
 -\frac{(\lambda-\gamma)^2}{\gamma}
       \mathbb E[w_m(Y)Y^2].
 \label{eq:dc-hierarchy}
\end{equation}
For $\lambda=\gamma$, this increases to the attained compatible-law minimum
$\sigma=\omega^2/\gamma$, since
\[
 \mathbb EY^{2m+2}
 =\frac{\binom{2m+2}{m+1}}{4^{m+1}}=O(m^{-1/2}).
\]
A grid point with $|\lambda-\gamma|\le\eta$ loses at most
$(m+1)\eta^2/\gamma$.  A deterministic finite grid with simultaneous
coverage therefore permits choosing the largest certified lower bound
without knowing $\gamma$.  Independently, $u/a$ converges to $\gamma$ as
the moment error, exposure error divided by $h$, and $h$ vanish with
$C(0)<1$.

For clarity, a conservative sufficient sequence makes the statement
constructive and removes the need for an oracle choice of $\lambda$.  Use
raw target and reference readouts bounded by one, as in the binary detector
example.  Fix positive parameter envelopes for the circle and valid resulting generator
caps.  Suppose a known integer $B\ge0$ obeys
$\kappa\ge\kappa_{\min}\ge2^{-B}$.  In a fixed time unit, at stage $k$ take
\begin{equation}
 \begin{gathered}
 m=k,\qquad h_k=2^{-k},\qquad \tau_k\le2^{-4k},\\
 K_k=2k+2,\qquad n_k=n_{{\rm ref},k}=2^{(8B+16)k}.
 \end{gathered}
 \label{eq:dc-consistency-schedule}
\end{equation}
Use the finite grid $[0,k]\cap2^{-k}\mathbb Z$ for $\lambda$.  References
use a known coarse shutter bracket shrinking at the stated exposure scale,
so their upper exposure bound also has that order.  Allocate a fixed total
error probability over all covered means and grid scores, or a summable
stagewise allocation.  Eventually the grid contains a point within $2^{-k}$
of the fixed true $\gamma$.

The raw score ranges are at most a polynomial in $k$ times $2^{2k}$.
Correcting degree at most $2k+2$ for gain multiplies interval widths by at
most $2^{2B(k+1)}$.  A bounded-mean radius with the declared sample size is
therefore at most a polynomial in $k$ times $2^{-(2B+6)k}$; the logarithm
of the finite-family size is $O(k)$.  Gain uncertainty has the same
vanishing order up to fixed and polynomial factors, by differentiating
$\kappa^{-d}$ on the known positive gain interval.  The calibration
reference count is included in~\eqref{eq:dc-consistency-schedule}.

The generic exposure bound gives
$O(k^4\tau_k/h_k^2)=O(k^4 2^{-2k})$ for raw-score blur.  The bias-correction
blur is smaller.  Uniformly over the finite grid, the finite-step penalty
is $O(k^2 2^{-k})$.  The nearest grid point loses $O(k4^{-k})$ in
\eqref{eq:dc-hierarchy}, while the gamma upper bound converges to the true
$\gamma$.  Direct uniform concentration around the covered means now gives
\begin{equation}
 \begin{gathered}
 \max_{\lambda\text{ on the grid}}\widehat L_k
       \longrightarrow\sigma\quad\text{in probability},\\
 \sigma-\max\widehat L_k
       =O_{\mathbb P}(k^{-1/2})
       =O_{\mathbb P}((\log n_k)^{-1/2}).
 \end{gathered}
 \label{eq:dc-consistency}
\end{equation}
In the last expression the difference is signed; the stochastic error is
smaller than the displayed order.  This rate is a conservative sufficient
rate, not a minimax assertion.  Direct concentration proves convergence
even at a fixed advertised confidence level; coverage alone would not.

Numerical enclosure error must also vanish after degree and time-step
amplification.  For example, carrying certified elementary-operation
enclosures to $k^2+(4B+12)k$ binary digits is ample for the finite algebraic
construction just described.  This includes refining any logarithm-series
truncation and square-root enclosure, rather than only rounding a fixed
approximation to more digits.  An implementation that retains a fixed
absolute rounding floor remains a valid finite-stage certificate, but its
unchanged output is not asserted to satisfy~\eqref{eq:dc-consistency}.

This is an existence construction with increasing physical access, not a
practical acquisition recommendation.  It requires shrinking exposure,
increasing replica access and reference measurements.  Fixed positive
exposure is not removed by this finite-difference argument.  Reset and
reference energies remain separate acquisition costs.

\section{Sharpness with an unobserved stochastic drive}

The sharp hierarchy is not confined to constant angular drift.  Let $Z$ be
any smooth reversible uniformly elliptic diffusion on a compact connected
manifold without boundary, with generator $L_Z$ and positive smooth
stationary density $\pi_Z$.  For a smooth real function $c$ on that manifold,
consider the joint generator
\begin{equation}
 L=\gamma\partial_\phi^2+c(z)\partial_\phi+L_Z,
 \qquad \phi\in\mathbb S^1,\qquad \gamma>0.
 \label{eq:dc-modulated-generator}
\end{equation}
All coordinates are subject to ordinary, even-variable reversal.  The
stationary measure is $(d\phi/2\pi)\pi_Z(dz)$: the angular drift is
divergence-free with respect to the uniform angular measure, and $L_Z$
preserves $\pi_Z$.  The diffusion tensor is block diagonal with angular
entry $\gamma$.  Thus
\begin{equation}
 S=\gamma\partial_\phi^2+L_Z,\qquad
 A=c(z)\partial_\phi,\qquad
 \sigma=\frac{\mathbb E_{\pi_Z}c^2}{\gamma}.
 \label{eq:dc-modulated-entropy}
\end{equation}

\begin{proposition}[Sharp minimum under reversible hidden modulation]
For $f(\phi,z)=\cos\phi$, the full stationary scalar observation law is
time-reversible.  Nevertheless its minimum compatible entropy over the
ordinary compact hidden-diffusion class of the observed-geometry theorem is
exactly $\mathbb E_{\pi_Z}c^2/\gamma$.  The actual joint diffusion attains
that value.  The finite polynomial hierarchy has the same sharp limit and
the same ideal formula~\eqref{eq:dc-hierarchy}, with
$\omega^2$ replaced by $\mathbb E_{\pi_Z}c^2$.
\end{proposition}
\begin{proof}
Direct computation gives
\[
 \Gamma_f=\gamma(1-f^2),\qquad Sf=-\gamma f,\qquad
 Af=-c(z)\sin\phi.
\]
At stationarity $Z$ and $\phi$ are independent, and
$U=(f,\Gamma_f)$ contains only $Y=f$.  The conditional quantities in the
population observed-geometry theorem therefore satisfy
\begin{align*}
 g(Y)&=-\gamma Y,\\
 K(Y)&=\gamma^2Y^2
       -\mathbb E_{\pi_Z}c^2\,(1-Y^2).
\end{align*}
The integrated conditional bound is
\[
 \mathbb E\left[
 \frac{[g(Y)^2-K(Y)]_+}{\gamma(1-Y^2)}
       \boldsymbol 1_{\{|Y|<1\}}\right]
 =\frac{\mathbb E_{\pi_Z}c^2}{\gamma}.
\]
The excluded endpoints have zero stationary probability.  This functional is
fixed by the ideal continuous observation law, so it bounds every compatible
ordinary hidden diffusion from below.  Equation~\eqref{eq:dc-modulated-entropy}
provides a realization attaining it.

For reversibility of the observation, the reflection
$\mathcal R(\phi,z)=(-\phi,z)$ preserves the stationary measure, fixes $f$
and conjugates $L$ to $L^\dagger$.  Applying this reflection to the reversed
stationary hidden path thus restores its forward law without changing any
detector values.  All finite-dimensional distributions of $Y$ agree with
their time reverses.  Finally, stationarity and independence yield
\begin{align*}
 \frac1\gamma\mathbb E\left[
 w_m(Y)\{(Af)^2-(-\lambda Y-Sf)^2\}\right]
 &={}\frac{\mathbb E c^2}{\gamma}
       (1-\mathbb EY^{2m+2})\\
 &\quad-\frac{(\lambda-\gamma)^2}{\gamma}
       \mathbb E[w_m(Y)Y^2],
\end{align*}
which proves the hierarchy assertion.
\end{proof}

The same sufficient acquisition and precision schedule applies under valid
uniform generator and positive-gain envelopes: the stationary arcsine
marginal and ideal hierarchy formula are unchanged, while the finite-lag
and exposure bounds already hold for the general hidden class.

The generator calibration can also be written explicitly.  Let
$C=\|c\|_\infty$ and $G=\|L_Zc\|_\infty$.  Then
\begin{align}
 Lf&=-\gamma\cos\phi-c\sin\phi,\\
 L^2f&=(\gamma^2-c^2)\cos\phi
                  +(2\gamma c-L_Zc)\sin\phi.
 \label{eq:dc-modulated-caps}
\end{align}
The adjoint formulas reverse the sine terms.  Both first-order norms equal
$\sqrt{\gamma^2+C^2}$, and both second-order norms equal
\[
 \sup_z\sqrt{(\gamma^2-c(z)^2)^2+
                       (2\gamma c(z)-L_Zc(z))^2}.
\]
The simple choices $M_{1,\pm}=\gamma+C$ and
$M_{2,\pm}=(\gamma+C)^2+G$ consequently suffice.  For unknown parameters,
calibrated upper envelopes for these expressions are required, just as in
the main finite certificate; the records do not automatically certify them.

For an explicit zero-mean example take a second angle $z$,
$L_Z=\eta\partial_z^2$ with $\eta>0$, and $c(z)=\varepsilon\cos z$.
Then
\begin{equation}
 \begin{gathered}
 \mathbb E c=0,\qquad
 \sigma=\frac{\varepsilon^2}{2\gamma},\qquad
 M_{1,\pm}=\sqrt{\gamma^2+\varepsilon^2},\\
 M_{2,\pm}=\sqrt{(\gamma^2-\varepsilon^2)^2
                            +(2\gamma+\eta)^2\varepsilon^2}.
 \end{gathered}
 \label{eq:dc-zero-mean-drive}
\end{equation}
The last identity follows by maximizing the squared norm over
$c^2\in[0,\varepsilon^2]$; its derivative in $c^2$ is positive.  Mean angular
current vanishes, but nonzero conditional angular currents remain and the
ordinary-reversal entropy minimum is positive when $\varepsilon\ne0$.
This is an analytically verified family, not an additional simulation or
physical experiment.  Interpreting a hidden coordinate as odd under
physical reversal changes the admissible class and can change the entropy
target; that interpretation is not included here.

\section{Scope and failure modes}

The theorem calibrates a scale in a declared pointwise geometry.  It does not
learn arbitrary hidden diffusivity or global generator suprema from a finite
record.  Rare-state countermodels obstruct that stronger inference.  A
false structural identity, correlated detector randomization, an altered
reference response, nonstationary target records or uncharged trial
dependence invalidates the certificate.  A single noisy center value cannot
replace the distinct replicas.

Replica inference, confidence-region inversion and the semigroup and
Dirichlet identities used here are established tools.  The result assembled
here is their explicit finite-sample propagation into this entropy
functional and its polynomial limiting hierarchy.  The record count,
replica channels, references, exposure and evolution time are acquisition
resources.  Total energy requires a separate calibrated sensor and preparation
model; it is not inferred from the sample count.
